% =====================================================================
% Robotic Arm Design & 3D Printing  +  Computer Vision Subsystem
% (~430 words combined; intentionally lighter on assembly per author)
% =====================================================================
\section{Mechanical Design and Computer-Vision Subsystem}
\label{sec:cad}

\subsection{Mechanical design and 3D printing}
The mechanical structure is a 4-DOF articulated arm with a single-DOF
gripper, built from PLA 3D-printed parts and standard hobby servos.
The link geometry is derived from the EEZYbotARM family with the
following original modifications: a redesigned base plate to mount the
HC-05, microcontroller and a power distribution board on a single
removable tray; a custom mounting bracket for the joint potentiometer
on each axis; and a bespoke gripper with a wider opening to accommodate
the coloured blocks used in the autonomous test set.

% TODO FIGURE: SolidWorks render or screenshot of the arm assembly.
\begin{figure}[t]
  \centering
  \fbox{\parbox{0.9\columnwidth}{\centering\vspace{4pt}
  \emph{[CAD render of the arm assembly --- SolidWorks isometric view.
  To be exported as PDF/PNG and included here.]}\vspace{4pt}}}
  \caption{Isometric CAD render of the arm assembly (placeholder).}
  \label{fig:cadrender}
\end{figure}

Parts were printed on an FDM printer using PLA at 0.2\,mm layer
height, 30\,\% gyroid infill, 3 perimeters and tree supports where
overhangs exceeded 60$^\circ$. Material was chosen for its dimensional
stability and the absence of post-processing requirements at this
load class; tougher materials such as PETG or carbon-fibre nylon
were rejected on the basis that the arm is bench-only and never
subject to thermal cycling. Iterative DFM tweaks made between the
first and second prints included widening the M3 servo-horn pockets
by 0.2\,mm to absorb FDM tolerance, adding a chamfer to the elbow
linkage to clear the shoulder bracket at full extension, and
relocating the potentiometer mounts to keep the wiper wires inside
the link rather than on the outside surface.

\subsection{Computer-vision subsystem}
\label{sec:vision}
The autonomous mode requires the arm to identify the colour of the
object presented to it. Two iterations of the vision pipeline were
built.

\subsubsection{First iteration --- ESP32 WROVER CAM}
A Freenove ESP32 WROVER CAM was initially used. The camera
streams QQVGA RGB565 frames into the ESP32, which converts each pixel
to integer HSV and bins it into one of four colour classes (red,
green, blue, yellow) with a 4\,\% presence threshold to suppress
background. The result is published as a one-line ASCII record over
\texttt{Serial2} to the Mega's \texttt{Serial1}:
\begin{center}\ttfamily\small
CAM,<state>,<colour>,<pct>$\backslash$n
\end{center}
This subsystem was wired in over a single GPIO~17~$\rightarrow$~Mega
pin~19 line plus shared ground; no level shift is required because
the ESP32 transmits 3.3\,V into the Mega's 5\,V tolerant RX. Initial
testing confirmed correct classification of the four colours under
controlled bench lighting.

% TODO PHOTO: ESP32 + camera mounted on bench, with the laptop
% showing the live colour-classification output.
\begin{figure}[t]
  \centering
  \fbox{\parbox{0.9\columnwidth}{\centering\vspace{4pt}
  \emph{[Photograph of the ESP32 WROVER CAM mounted to the arm base
  during the first iteration, with the live colour classification
  visible in the serial monitor on the laptop.]}\vspace{4pt}}}
  \caption{ESP32 WROVER CAM during the first vision iteration
  (photo placeholder).}
  \label{fig:esp32cam}
\end{figure}

\subsubsection{Second iteration --- Raspberry Pi camera module}
In extended testing the ESP32 pipeline proved sensitive to ambient
lighting changes and exhibited occasional false classifications at
the boundary between green and yellow. The pipeline was therefore
re-implemented on a Raspberry Pi with a CSI camera module, where a
larger frame size, floating-point HSV conversion and a small
classical-CV pipeline (Gaussian blur, HSV thresholding,
contour-area filtering) gave substantially more reliable
classification across a wider range of conditions. The Pi reports
its result to the Mega using the same one-line ASCII protocol on
the same Mega UART pin, so no firmware changes were required on
the arm side --- the Mega is agnostic to which co-processor is
producing the frames.

% TODO PHOTO: Pi + camera module mounted to arm.
\begin{figure}[t]
  \centering
  \fbox{\parbox{0.9\columnwidth}{\centering\vspace{4pt}
  \emph{[Photograph of the Raspberry Pi + camera module mounted to
  the arm during the second vision iteration. To be added once the
  Pi pipeline is fully integrated.]}\vspace{4pt}}}
  \caption{Raspberry Pi + camera module during the second iteration
  (photo placeholder).}
  \label{fig:picam}
\end{figure}

The Pi is treated strictly as a vision co-processor: all real-time
motor control remains on the Mega, preserving the embedded-systems
character of the project. This is consistent with the Pi vs Arduino
position summarised in Section~\ref{sec:literature}.

\corehit{Brief Objective~3 and the CAD/3D-printing rubric: original
CAD modifications, FDM print parameters justified, fabrication
performed in-house.}

\advhit{Two iterations of an embedded computer-vision pipeline
(on-device HSV on ESP32; classical CV on the Pi); shared
single-protocol UART link so the arm-side firmware is unchanged
between iterations; multi-MCU heterogeneous architecture with
clean separation of real-time and non-real-time concerns.}
